\section{Related Work}
\textbf{Disaggregated Memory and Prefix Caching.}
Prior work decouples KV cache management from compute to mitigate fragmentation and imbalance. TokenLake~\cite{tokenlake}, for example, proposes a disaggregated memory architecture with unified prefix pools to reduce redundancy across instances. While effective at improving spatial efficiency, such approaches do not address the \emph{temporal} overcommitment inherent in offline agentic inference, where concurrent agents jointly expand context during the middle phase and overwhelm even a unified pool. Our work complements these designs by introducing proactive flow control to prevent memory saturation.

\textbf{Application-Centric Serving.}
Recent systems optimize beyond request-level scheduling by exploiting application structure to reduce latency and job completion time. Parrot~\cite{parrot}, Teola~\cite{teola}, Autellix~\cite{autellix}, and Kairos~\cite{kairos} leverage DAGs, dependency awareness, or memory prediction to reorder requests and mitigate blocking in multi-turn workloads. They assume that memory pressure can be handled via eviction or reordering. In contrast, our large-batch offline setting requires explicit \emph{admission control}, as excessive concurrency induces thrashing regardless of scheduling order.

\textbf{Agent-Native Inference and Concurrency.}
Agent execution interleaves GPU generation with CPU-side tool calls, complicating memory management. TokenCake~\cite{tokencake} and Continuum~\cite{continuum} address this via reactive mechanisms such as offloading and TTL-based cache pinning. Our approach is fundamentally different: inspired by network congestion control, we proactively regulate agent concurrency using AIMD to bound memory pressure and eliminate middle-phase thrashing, avoiding reliance on costly swap operations.